\sectionkicker{Source-backed technical notes}
\subsection*{強さだけでは測れない――Foundation Models、TruthfulQA、リスク分類: Technical Notes}
本欄は記事本文で圧縮した一次資料上の情報を、比較・再検証しやすい形で整理したものである。時系列、確認済みの主張、留意点、一次資料URLを掲載する。
\smallskip
\noindent{\small\textit{Technical Notesでは、一次資料から確認した公開・提供・研究上の事実を記載する。数値・能力評価や提供元・著者による比較表現は、独立再現済みの結果ではなく帰属claimとして扱う。}}\par
\medskip
\subsection*{Theme at a glance}
\addcontentsline{toc}{subsection}{Theme at a glance}
\begin{center}
\footnotesize
\begin{tabularx}{\linewidth}{@{}p{0.20\linewidth}p{0.10\linewidth}p{0.14\linewidth}X@{}}
\toprule
資料 & 位置づけ & 種別 & 時系列 \\
\midrule
Foundation Models report & 主要資料 & 論文 & 2021-08-16 \\
TruthfulQA & 主要資料 & 評価ベンチマーク & 2021-09-08 \\
Ethical and social risks of harm from Language Models & 補足資料 & 論文 & 2021-12-08 \\
Training Verifiers / GSM8K & 補足資料 & 論文 & 2021-10-27, 2021-10-29 \\
\bottomrule
\end{tabularx}
\end{center}
\normalsize

\begin{technicalnote}{Foundation Models report}{主要資料}
\begingroup
% reader-facing Technical Notes break policy
\widowpenalty=10000
\clubpenalty=10000
\displaywidowpenalty=10000
\begin{tabularx}{\linewidth}{@{}>{\bfseries}p{0.22\linewidth}X@{}}
組織 & Stanford CRFM / authors \\
種別 & 論文 \\
時系列 & 2021-08-16 \\
\end{tabularx}
\smallskip
{\bfseries 一次資料から整理したtechnical points}
\begin{itemize}[leftmargin=1.5em,itemsep=0.35em]
\item \textbf{一次情報で確認できる事実}: Stanford CRFM / authorsの選定済み一次資料では、Foundation Models reportはbroad data at scaleでtrainingされ、多数のdownstream taskへadaptされるmodel群を`foundation models`として整理した。 同じ基盤が多用途へ再利用されるhomogenizationの利点と、基盤側のdefectがdownstreamへ継承されるriskを同じframeworkで論じた。
\end{itemize}
\begin{minipage}{\linewidth}
% reader-facing Technical Notes generic boundary/source tail group
\begin{samepage}
{\bfseries 一次資料}
\begingroup\sloppy
\Urlmuskip=0mu plus 2mu\relax
\begin{itemize}[leftmargin=1.5em,itemsep=0.25em]
\item {\scriptsize\url{http://arxiv.org/abs/2108.07258v3}}
\end{itemize}
\endgroup
\end{samepage}
\end{minipage}
\endgroup
\end{technicalnote}
\medskip

\begin{technicalnote}{TruthfulQA}{主要資料}
\begingroup
% reader-facing Technical Notes break policy
\widowpenalty=10000
\clubpenalty=10000
\displaywidowpenalty=10000
\begin{tabularx}{\linewidth}{@{}>{\bfseries}p{0.22\linewidth}X@{}}
組織 & OpenAI / authors \\
種別 & 評価ベンチマーク \\
時系列 & 2021-09-08 \\
\end{tabularx}
\smallskip
{\bfseries 一次資料から整理したtechnical points}
\begin{itemize}[leftmargin=1.5em,itemsep=0.35em]
\item \textbf{一次情報で確認できる事実}: OpenAI / authorsの選定済み一次資料では、TruthfulQAはhuman misconceptionに引きずられずtruthfulな回答を生成できるかを測るため、38 category・817 questionからなるbenchmarkを構成した。 GPT-3、GPT-Neo/J、GPT-2、T5系modelを比較し、model scaleだけではtruthfulnessが改善しないケースを評価対象にした。
\end{itemize}
\begin{minipage}{\linewidth}
% reader-facing Technical Notes generic boundary/source tail group
\begin{samepage}
{\bfseries 一次資料}
\begingroup\sloppy
\Urlmuskip=0mu plus 2mu\relax
\begin{itemize}[leftmargin=1.5em,itemsep=0.25em]
\item {\scriptsize\url{http://arxiv.org/abs/2109.07958v2}}
\item {\scriptsize\url{https://openai.com/index/truthfulqa}}
\end{itemize}
\endgroup
\end{samepage}
\end{minipage}
\endgroup
\end{technicalnote}
\medskip

\begin{technicalnote}{Ethical and social risks of harm from Language Models}{補足資料}
\begingroup
% reader-facing Technical Notes break policy
\widowpenalty=10000
\clubpenalty=10000
\displaywidowpenalty=10000
\begin{tabularx}{\linewidth}{@{}>{\bfseries}p{0.22\linewidth}X@{}}
組織 & DeepMind / authors \\
種別 & 論文 \\
時系列 & 2021-12-08 \\
\end{tabularx}
\smallskip
{\bfseries 一次資料から整理したtechnical points}
\begin{itemize}[leftmargin=1.5em,itemsep=0.35em]
\item \textbf{一次情報で確認できる事実}: DeepMind / authorsの選定済み一次資料では、論文はlarge-scale language modelのrisk landscapeを6つのrisk areaに整理し、discrimination/toxicity、information hazards、misinformation、malicious use、human-computer interaction、automation/access/environmental harmsを分析した。 合計21 riskを掘り下げ、risk origin、mitigation approach、組織的責任とparticipationまでを同じtaxonomyで扱った。
\end{itemize}
\begin{minipage}{\linewidth}
% reader-facing Technical Notes generic boundary/source tail group
\begin{samepage}
{\bfseries 一次資料}
\begingroup\sloppy
\Urlmuskip=0mu plus 2mu\relax
\begin{itemize}[leftmargin=1.5em,itemsep=0.25em]
\item {\scriptsize\url{http://arxiv.org/abs/2112.04359v1}}
\end{itemize}
\endgroup
\end{samepage}
\end{minipage}
\endgroup
\end{technicalnote}
\medskip

\begin{technicalnote}{Training Verifiers / GSM8K}{補足資料}
\begingroup
% reader-facing Technical Notes break policy
\widowpenalty=10000
\clubpenalty=10000
\displaywidowpenalty=10000
\begin{tabularx}{\linewidth}{@{}>{\bfseries}p{0.22\linewidth}X@{}}
組織 & OpenAI / authors \\
種別 & 論文 \\
時系列 & 2021-10-27, 2021-10-29 \\
\end{tabularx}
\smallskip
{\bfseries 一次資料から整理したtechnical points}
\begin{itemize}[leftmargin=1.5em,itemsep=0.35em]
\item \textbf{一次情報で確認できる事実}: OpenAI / authorsの選定済み一次資料では、GSM8Kは8.5K件のgrade-school math word problemを収録し、多段の数学推論を評価するdatasetとして提示された。 手法側ではmodelから複数のcandidate solutionを生成し、completionの正しさを判定するverifierでrankingして最上位を選ぶ構成を評価した。
\end{itemize}
\begin{minipage}{\linewidth}
% reader-facing Technical Notes generic boundary/source tail group
\begin{samepage}
{\bfseries 一次資料}
\begingroup\sloppy
\Urlmuskip=0mu plus 2mu\relax
\begin{itemize}[leftmargin=1.5em,itemsep=0.25em]
\item {\scriptsize\url{http://arxiv.org/abs/2110.14168v2}}
\item {\scriptsize\url{https://openai.com/index/solving-math-word-problems}}
\end{itemize}
\endgroup
\end{samepage}
\end{minipage}
\endgroup
\end{technicalnote}
\medskip
